% Shared colophon: how the edition was produced + the project's goals.
% Input by each edition's about.tex, inside a centred minipage.

\medskip
{\booktitlefont\scshape How this edition was made}\par
This volume is typeset with LuaLaTeX and David Purton's \texttt{scripture}
package, set in Georg Duffner's EB~Garamond on a strict baseline grid. The
scripture text is drawn from its publisher's source, converted to
scripture-package markup by the project's Python generators, and composed in the
manner of the early printed Bibles: hanging (margin-protruding) verse numbers,
drop-cap lettrine chapter openings, the divine name in small capitals, and
period public-domain woodcut headpieces vectorized from early books and assigned
by canonical division. Proper names are hyphenated after a published
Bible-typesetting authority, and widows and orphans are resolved without
disturbing the baseline grid.

\medskip
{\booktitlefont\scshape About the project}\par
\emph{geneve\_1564} is an open typesetting project in the visual tradition of the
Geneva Bible of 1564. Its goal is to render Scripture as a beautiful, readable
book --- honouring historic printing practice --- that anyone can build for
themselves, for print or for e-ink readers. It began as a fork of Rapha\"el
Pinson's Geneva~1564 reproduction and was inspired by the \texttt{scripture}
package. Source, provenance, and build instructions:
\texttt{github.com/coopermj/geneve\_1564}.
